% Generated by scripts/build_paper_tables.py; do not edit by hand.
\par
\begin{table}[htbp]
\centering
\caption{按月选参与固定参数比较（2025年4—12月）}\label{tab:rolling-summary}
\footnotesize
\setlength{\tabcolsep}{4pt}
\renewcommand{\arraystretch}{1.15}
\begin{tabular}{lrrrrr}
\toprule
方案 & 参数规则 & 总费用（元） & 紧急电（kWh） & 期初SOC & 期末SOC \\
\midrule
问题3 & 固定1月参数 & 11,025,848.81 & 24,392.62 & 6,461.68 & 5,987.72 \\
问题3 & 上月选择融合 & 11,027,957.95 & 39,352.86 & 6,461.68 & 5,987.72 \\
问题3 & 上月选择纯历史 & 11,132,033.28 & 37,016.85 & 6,461.68 & 5,987.72 \\
问题4-3 & 固定1月参数 & 11,695,213.43 & 23,070.56 & 6,517.67 & 5,997.11 \\
问题4-3 & 上月选择融合 & 11,711,534.09 & 39,935.96 & 6,517.67 & 5,997.11 \\
问题4-3 & 上月选择纯历史 & 11,812,014.14 & 46,519.39 & 6,517.67 & 5,997.11 \\
\bottomrule
\end{tabular}
\par\smallskip
\begin{minipage}{0.97\linewidth}\footnotesize 三条路径共用4月1日期初库存，此后逐月连续；SOC单位为kWh。本表覆盖275天，不与2—12月334天总费直接相减。\end{minipage}
\end{table}
\par
